\documentclass[acmsmall,nonacm,screen]{acmart}

\usepackage{mathpartir}
\usepackage{booktabs}

\citestyle{acmauthoryear}
\settopmatter{printacmref=false}

\AtEndPreamble{\theoremstyle{acmdefinition}\newtheorem{remark}[theorem]{Remark}}

\newcommand{\kw}[1]{\mathsf{#1}}
\newcommand{\nm}[1]{\mathit{#1}}
\newcommand{\letin}[3]{\kw{let}\;#1 = #2\;\kw{in}\;#3}
\newcommand{\ite}[3]{\kw{if}\;#1\;\kw{then}\;#2\;\kw{else}\;#3}
\newcommand{\iter}[4]{\kw{iter}\;#1\;\kw{from}\;#2\;\kw{as}\;#3.\,#4}
\newcommand{\rd}[1]{\kw{read}\;#1}
\newcommand{\wrt}[2]{\kw{write}\;#1\;#2}
\newcommand{\mdl}[1]{\kw{model}\;#1}
\newcommand{\chk}[1]{\kw{check}\;#1}
\newcommand{\scan}{\kw{scan}}
\newcommand{\sem}[1]{[\![#1]\!]}
\newcommand{\evto}{\Downarrow}
\newcommand{\fp}[1]{\mathrm{fp}(#1)}
\newcommand{\dom}{\mathrm{dom}\,}
\newcommand{\upd}{\oplus}

% A display for programs. It measures its content and warns when the content is wider than the line.
\newsavebox{\progbox}
\newenvironment{prog}
  {\begin{lrbox}{\progbox}$\displaystyle\begin{array}{@{}l@{\qquad}l@{}}}
  {\end{array}$\end{lrbox}%
   \ifdim\wd\progbox>\linewidth
     \PackageWarning{paper}{Overfull program display by \the\dimexpr\wd\progbox-\linewidth\relax}%
   \fi
   \[\usebox{\progbox}\]\ignorespacesafterend}

% written by merge.py; do not edit
\newcommand{\exCompared}{216}
\newcommand{\exConeFileFleet}{5}
\newcommand{\exConeFileHalf}{337}
\newcommand{\exConeFileMax}{2,411}
\newcommand{\exConeFileMean}{456.1}
\newcommand{\exConeFileMedian}{91}
\newcommand{\exConeFileP}{18.88}
\newcommand{\exConePrivateFleet}{15,733}
\newcommand{\exConePrivateMax}{24}
\newcommand{\exConePrivateMean}{0.2}
\newcommand{\exConePrivateMedian}{0}
\newcommand{\exConePrivateNone}{95}
\newcommand{\exConePrivateP}{0.01}
\newcommand{\exConePublicFleet}{8}
\newcommand{\exConePublicHalf}{228}
\newcommand{\exConePublicMax}{2,410}
\newcommand{\exConePublicMean}{299.6}
\newcommand{\exConePublicMedian}{32}
\newcommand{\exConePublicP}{12.40}
\newcommand{\exCyclic}{614}
\newcommand{\exEntries}{33,676}
\newcommand{\exFleetArrivea}{1.0}
\newcommand{\exFleetArriveb}{4.1}
\newcommand{\exFleetArrivec}{14.4}
\newcommand{\exFleetArrived}{45.0}
\newcommand{\exFleetArrivee}{120.3}
\newcommand{\exFleetBesta}{0.12}
\newcommand{\exFleetBestb}{0.47}
\newcommand{\exFleetBestc}{2.31}
\newcommand{\exFleetBestd}{11.53}
\newcommand{\exFleetBeste}{53.32}
\newcommand{\exFleetCyclica}{11}
\newcommand{\exFleetCyclicb}{39}
\newcommand{\exFleetCyclicc}{94}
\newcommand{\exFleetCyclicd}{100}
\newcommand{\exFleetCyclice}{100}
\newcommand{\exFleetHota}{0.51}
\newcommand{\exFleetHotb}{1.64}
\newcommand{\exFleetHotc}{5.12}
\newcommand{\exFleetKpa}{0.25}
\newcommand{\exFleetKpb}{0.50}
\newcommand{\exFleetKpc}{0.99}
\newcommand{\exFleetKpd}{1.98}
\newcommand{\exFleetKpe}{3.96}
\newcommand{\exFleetMutual}{0.31}
\newcommand{\exFleetP}{2.5}
\newcommand{\exFleetRatio}{8}
\newcommand{\exFleetSizeArrive}{13}
\newcommand{\exFleetSizeEdges}{36}
\newcommand{\exFree}{2,386}
\newcommand{\exImportsAll}{283}
\newcommand{\exImportsPlain}{3,725}
\newcommand{\exImportsPublic}{6,357}
\newcommand{\exModules}{2,416}
\newcommand{\exQueueCalls}{1,160}
\newcommand{\exQueueSaved}{6,559}
\newcommand{\exQueueTasks}{9,071}
\newcommand{\exQueues}{3,000}
\newcommand{\exRebaseKept}{2,458}
\newcommand{\exRebaseMade}{975}
\newcommand{\exRebaseResumed}{2,035}
\newcommand{\exRebases}{4,000}
\newcommand{\exRepeated}{2,868}
\newcommand{\exRepeatedPercent}{9}
\newcommand{\exResumed}{2,400}
\newcommand{\exResumedKeptPercent}{37}
\newcommand{\exSavedPercent}{85}
\newcommand{\exStepLastDelta}{2}
\newcommand{\exStepScanLocal}{2}
\newcommand{\exStepSkipCheck}{22}
\newcommand{\exStepSkipRead}{9}
\newcommand{\exStepStaleDelta}{782}
\newcommand{\exSteps}{5}
\newcommand{\exUntouched}{6,455}


\begin{document}

\title{A Merge Is a Schedule: What Parallel Agents Must Redo}

\author{Claude-Fable-5.2}
\authornote{An AI model. The proofs are on paper and have not been checked by a machine. A
script implements the calculus and tests the theorems on random tasks; a second script reads
the import graph of a library. At the time of writing nobody has reviewed the paper.}

\begin{abstract}
Agents now work in fleets. Each takes a copy of a repository, works for an hour, and proposes
a change, and version control merges the changes that touch different lines. Whether the
merged state is one that any of the agents would have produced, or would have accepted, is
another matter: each of them read files and ran tests on a state that the others have changed
in the meantime. We treat the work of an agent as a transaction on a snapshot. Its journal
records what it read, what it asked the model, what it checked, and what it wrote, and every
answer carries a footprint, the keys of the repository on which it depends. To merge is then
to replay each journal on the state that the earlier ones leave. A request whose footprint no
earlier task has changed keeps its answer. Another one is made again and compared. The model
is not asked, until an answer differs; from there the task goes on live. We prove that the
result is the state of running the tasks one after the other, with the answers that the model
gave; that tasks which have not read each other's writes merge with no work at all, in an
order that the footprints determine; that with a cycle of such reads someone has to resume,
and that the union of the writes can then be a state that no order produces; and that a check
travels with its footprint, so that only the checks whose footprints were changed are run
again. A script tests the theorems on random tasks: in \exQueues{} queues, \exSavedPercent{}
percent of the model's answers are kept and \exRepeatedPercent{} percent of the reads and
checks are made again. For a fleet with random footprints, a queue in the order of arrival
makes a number of tasks resume that grows with the square of the fleet, and an order along
the footprints makes it grow about \exFleetRatio{} times slower. In Lean's own library a change of a
proof makes \exConePrivateMean{} other modules be checked again on average, and a change of a
statement \exConePublicMean{}.
\end{abstract}

\keywords{merging, snapshot isolation, serializability, incremental checking, replay, agents}

\authorsaddresses{}

\maketitle

\section{Introduction}

A coding agent that works alone is slow, so people run many. Each agent gets a copy of the
repository, a worktree or a container, and a task. It reads files, asks its model, edits, and
runs the tests until they pass. Then it proposes its change. Version control merges the
changes of several agents if they do not touch the same lines.

That a merge has no conflict in this sense says little. Suppose one agent renames a function
and updates every caller that it finds, while another adds a new caller. The two changes
touch different lines, each agent saw its tests pass, and the merged program does not
compile. Studies of real histories find such merges regularly: accepted by version control,
then failing to build or failing tests~\citep{brun2011}. The usual remedy is to run all tests
again on the merged state, and, if they fail, to send one of the agents back to work on top
of the other's change. For a fleet both steps are expensive. The tests of a large repository
take long, and an agent that starts again pays for all its model calls a second time.

How much of this is needed? An agent's work is a function of what it read. If nothing that
it read has changed, the work is still what the agent would do now, and the tests that it ran
would give the same verdicts. So the questions are what an agent read, and what the others
wrote. Databases have asked these questions for decades. An agent on a copy of the repository
is a transaction under snapshot isolation~\citep{berenson1995}, and the merge of several is
correct in the sense that matters if it is \emph{serializable}: the same as running the tasks
one after the other in some order~\citep{papadimitriou1979}.

This paper works the correspondence out for agents, where two things are special. An agent's
reads are known exactly, because an agent that is run durably keeps a journal of every
request~\citep{tpp8}. And its steps differ in cost by orders of magnitude: a read costs
nothing, a test costs seconds, a call of the model costs money. So we do not ask only whether
a task has to be done again. We ask which of its steps do.

\paragraph{Replay.} A task runs on a base state and leaves a journal. Each entry is a request
with its answer, and each answer of the repository carries a \emph{footprint}, the keys of
the base state on which it depends. To move the task to another state, we replay the journal
there. An answer whose footprint lies outside the keys that have changed is kept. Another is
obtained again and compared. A question to the model is answered from the journal. If an
answer differs, the task \emph{resumes}: from this request on it runs live. We prove
(Theorem~\ref{thm:replay}) that the result is a run of the task on the new state, in which
the model gave the recorded answers up to the point of resumption. Rebasing a task costs the
requests that are made again, and the model calls after the first answer that differs.

\paragraph{The queue.} A merge queue replays the journals one after the other, each on the
state that the earlier ones leave. The final state is that of the tasks run one after the
other in the order of the queue (Theorem~\ref{thm:schedule}): a merge is a schedule. The
order matters for the cost and not for the correctness. Draw an edge from task $i$ to task
$j$ if $i$ read a key that $j$ changed. If the graph has no cycle, an order along its edges
replays every journal without making a single request (Corollary~\ref{cor:free}), and the
merged state is the union of the writes. In general the tasks that resume can be confined to
a feedback vertex set (Corollary~\ref{cor:feedback}). With a cycle the union of the writes
may be a state that no order produces, and it may break a check that every order keeps
(Proposition~\ref{prop:cycle}).

\paragraph{Checks.} A test is a request like any other: its answer has a footprint, the keys
that the test read. It is kept across a merge if no other task changed its footprint
(Corollary~\ref{cor:checks}). This is the sense in which changes can be merged without
rechecking, and it is the notion of a verifying trace in build systems~\citep{mokhov2018},
applied to merges.

\paragraph{Fleets and libraries.} Two sections put numbers on this. Section~\ref{sec:fleet}
takes a fleet of $k$ tasks with random footprints. In the order of arrival the expected
number of tasks that resume is about $p k^2 / 2$, where $p$ is the probability that one task
read what another wrote. In an order along the edges it is about $p_2 k^2 / 2$, where $p_2$
is the probability that two tasks read each other's writes; in our example $p$ is
\exFleetP{} percent and $p_2$ is \exFleetRatio{} times smaller. Section~\ref{sec:library} measures
footprints in a real library, the \exModules{} modules of Lean's own sources. If a check
reads everything that its module imports, a change of one module makes \exConeFileMean{}
modules be checked again on average. With the scopes of Lean's module
system~\citep{leanmodules}, a change of a statement makes \exConePublicMean{} modules be
checked again, and a change that stays inside proofs \exConePrivateMean{}.

\paragraph{Tests.} The proofs are on paper. An artifact implements tasks, replay, and the
queue, compares every queue with the tasks run one after the other, and breaks the replay in
\exSteps{} ways to show that each of its steps is needed (Section~\ref{sec:tests}).

\section{Two Tasks, One Repository}
\label{sec:overview}

\subsection{Each Keeps the Check, Both Break It}
\label{sec:skew}

A repository holds two switches, $a$ and $b$, both on. A check requires that at least one
switch is on. Two agents get the same instruction: turn your switch off if the other one is
on.
\begin{prog}
\text{task 1:}\quad \letin{x}{\rd{b}}{\ite{x = 1}{\wrt{a}{0}}{0}}\\
\text{task 2:}\quad \letin{x}{\rd{a}}{\ite{x = 1}{\wrt{b}{0}}{0}}
\end{prog}
Each task runs on its own copy. Task 1 reads $b$, finds it on, and turns $a$ off. Task 2
reads $a$, finds it on, and turns $b$ off. Each runs the check on its copy, and the check
passes. The two changes touch different keys, so version control merges them, and in the
merged state both switches are off.

No order of the two tasks ends there. If task 1 runs first, task 2 finds $a$ off and does
nothing, and the other way round. In databases this is called write
skew~\citep{berenson1995}. The renamed function of the introduction is the same pattern,
with the set of callers in the place of a switch.

The journals show what went wrong. The journal of task 1 is
$\rd{b} = 1 \cdot \wrt{a}{0}$, and that of task 2 is $\rd{a} = 1 \cdot \wrt{b}{0}$. Task 1
read what task 2 wrote, and task 2 read what task 1 wrote. A queue that replays the journals
merges task 1 first; nothing has changed yet, so the journal stands. Then it replays task 2
on the state in which $a$ is off. The first entry of the journal is a read of $a$, and $a$
has changed, so the read is made again. It now gives $0$, the recorded answer was $1$, and
the task resumes from there: it takes the other branch and writes nothing. The final state
has $a$ off and $b$ on. It is the state of running task 1 and then task 2, and the check
holds.

\subsection{One-Way Reads Are Free}

Not every stale read forces work. Give the two agents other tasks. Task 1 copies a
switch, and task 2 turns that switch off without looking at anything:
\begin{prog}
\text{task 1:}\quad \letin{x}{\rd{b}}{\wrt{c}{x}}\\
\text{task 2:}\quad \wrt{b}{0}
\end{prog}
Task 1 has read a key that task 2 changes, and task 2 has read nothing. If the queue takes
task 2 first, the replay of task 1 finds that $b$ has changed, reads $0$ where the journal
says $1$, and resumes; it writes $0$ to $c$. If the queue takes task 1 first, the journal of
task 1 stands, because nothing has changed yet, and then the journal of task 2 stands,
because it has no reads. No request is made.

Both orders are schedules, and they end in different states: $c$ is $0$ after the first and
$1$ after the second. Neither is wrong. The second is the one that the work of the agents
already describes, and it costs nothing. So a task that has read what another wrote goes
\emph{before} it. This is the rule of the paper in its smallest form: only a cycle of such
reads forces work.

\subsection{A Check Is a Read}

An agent that ran the tests has an entry $\chk{c} = 1$ in its journal. The test read some
keys while it ran, and they are the footprint of the entry. Here is the journal of a task
that reads $a$, asks the model what to write, writes $b$, and runs a test that reads $b$ and
$d$:
\[
\rd{a} = 1\;[a] \;\cdot\; \mdl{1} = 3 \;\cdot\; \wrt{b}{3} \;\cdot\; \chk{c} = 1\;[d].
\]
The footprint of the read is $a$. The footprint of the test is $d$ only: the test also read
$b$, but the task wrote $b$ itself, and the replay will write it again. Suppose another task,
merged earlier, has changed $d$ and nothing else. The replay keeps the read, since $a$ has
not changed. It keeps the answer of the model without asking, and it repeats the write. The
test is run again, because $d$ has changed. If it passes again, the journal stands, and all
that the merge has cost is one test. If it fails, the task resumes at this point, with a
failing test in front of it and the model's answer still in hand, which is where an agent
wants to be. Had the other task changed $a$, the read would have been compared first, and a
different value would have sent the task back to its first request.

A test whose footprint is the whole repository is run again after every merge. A linter
that looks at every file is of this kind, and so is a proof checker that consults a global
table. Section~\ref{sec:library} looks at how large the footprints of checks are in a real
library of proofs.

\section{The Calculus}
\label{sec:calculus}

\subsection{Tasks and Checks}

\begin{figure}[t]
\[
\begin{array}{@{}l@{\quad}r@{\;}l@{}}
\text{terms} & t ::= & x \mid n \mid t + t \\
\text{tests} & b ::= & t = t \mid t \le t \\
\text{expressions} & e ::= & t \mid \letin{x}{e}{e} \mid \ite{b}{e}{e} \mid \iter{n}{t}{x}{e} \\
 & \mid & \rd{t} \mid \wrt{t}{t} \mid \mdl{t} \mid \chk{c} \mid \scan \\[1ex]
\text{entries} & & \rd{k} = v\;[F] \mid \wrt{k}{v} \mid \mdl{n} = r \mid \chk{c} = v\;[F] \mid \scan = v\;[F]
\end{array}
\]
\Description{The grammar of terms, tests, expressions, and journal entries.}
\caption{Syntax. Numbers $n$, $v$, $r$ are values; $k$ ranges over the keys of the
repository, $c$ over the names of checks, and $F$ over sets of keys.}
\label{fig:syntax}
\end{figure}

A \emph{state} $\sigma$ of the repository gives a value to each of finitely many
\emph{keys}. A key stands for a file, or for a smaller unit such as a declaration.
Figure~\ref{fig:syntax} gives the expressions. A \emph{task} is a closed expression. It
computes with numbers, and it makes five kinds of requests: $\rd{t}$ returns the value of the
key that $t$ denotes; $\wrt{t}{t'}$ sets a key; $\mdl{t}$ asks the model, which may answer
with any number; $\chk{c}$ runs the check $c$ and returns its value; and $\scan$ returns a
value that depends on all keys, their sum. Which key a task reads may depend on what it has
read before and on what the model said, so the reads of a task are not known in advance.

A \emph{check} is a closed expression without $\kw{write}$, $\kw{model}$ and $\kw{check}$. It
is a deterministic function of the state, such as a test, a compiler run, or a proof checker.
A repository has a fixed family of checks, named $c$. The request $\scan$ models what looks
at everything: a directory listing, a global registry.

Loops are bounded, as in the companion calculi~\citep{tpp8}, so every task ends.

\subsection{Runs and Journals}
\label{sec:runs}

A task runs on a \emph{base} state $\sigma$ and keeps its writes to itself, in an
\emph{overlay} $\omega$, a partial state. What the task sees is its \emph{view}
$\sigma \upd \omega$, the base with the overlay on top. The judgment
\[
\sigma; \omega \vdash e \evto v; \omega'; J
\]
says that $e$, started with the overlay $\omega$, has the value $v$, ends with the overlay
$\omega'$, and leaves the journal $J$, a sequence of entries. The rules for terms, $\kw{let}$,
conditionals, and loops are the usual ones; they pass the overlay along and join the
journals. Figure~\ref{fig:runs} has the rules for the requests. A read looks at the view. If
the key comes from the base, the entry names it; if the task has written the key itself, the
entry depends on no key of the base. A write changes the overlay and nothing else. The
model may answer anything. A check evaluates its expression $q_c$ on the view, by a judgment
$\sigma; \omega \vdash q \Rightarrow v; F$ that has the same rules for reads and scans and
collects the keys that were read from the base. A scan depends on every key that the task
has not written.

\begin{figure}[t]
\begin{mathpar}
\inferrule*[right=Read]{v = (\sigma \upd \omega)(k) \\ F = \{k\} \setminus \dom \omega}
  {\sigma; \omega \vdash \rd{k} \evto v; \omega; \rd{k} = v\;[F]}
\and
\inferrule*[right=Write]{ }{\sigma; \omega \vdash \wrt{k}{n} \evto n; \omega[k \mapsto n]; \wrt{k}{n}}
\and
\inferrule*[right=Model]{ }{\sigma; \omega \vdash \mdl{n} \evto r; \omega; \mdl{n} = r}
\and
\inferrule*[right=Check]{\sigma; \omega \vdash q_c \Rightarrow v; F}
  {\sigma; \omega \vdash \chk{c} \evto v; \omega; \chk{c} = v\;[F]}
\and
\inferrule*[right=Scan]{v = \textstyle\sum_k (\sigma \upd \omega)(k) \\ F = \{k \mid k \notin \dom \omega\}}
  {\sigma; \omega \vdash \scan \evto v; \omega; \scan = v\;[F]}
\and
\inferrule*[right=Let]{\sigma; \omega \vdash e_1 \evto v_1; \omega_1; J_1 \\ \sigma; \omega_1 \vdash e_2[v_1/x] \evto v_2; \omega_2; J_2}
  {\sigma; \omega \vdash \letin{x}{e_1}{e_2} \evto v_2; \omega_2; J_1 J_2}
\end{mathpar}
\Description{Six rules of a big-step semantics with an overlay and a journal.}
\caption{Runs. The terms of a request have been evaluated: $k$ is a key and $n$ a number. In
\textsc{Model}, $r$ is any number.}
\label{fig:runs}
\end{figure}

The set $F$ is the \emph{footprint} of an answer. It is found while the answer is computed, so
it is exact for this state: a check that tests a key and then reads one of two others has
the two keys that it read, not three. The footprint $\fp{J}$ of a journal is the union of
the footprints of its entries. A \emph{run} of a task $e$ on $\sigma$ is a derivation of
$\sigma; \emptyset \vdash e \evto v; \omega; J$. A run is determined by the base state and
the answers of the model.

\subsection{Replay}
\label{sec:replay}

\begin{table}[t]
\caption{Replay of a journal on a state $\sigma$, given a set $\Delta$ that contains every
key on which $\sigma$ differs from the base of the journal. The task is evaluated as in a
run; this is what each request does while the replay has not resumed.}
\label{tab:replay}
\small
\begin{tabular}{@{}lll@{}}
\toprule
 & The entry of the journal & What the replay does \\
\midrule
1 & $\rd{k} = v\;[F]$, $\chk{c} = v\;[F]$ or $\scan = v\;[F]$, & keeps the entry; the answer is $v$ \\
  & with $F \cap \Delta = \emptyset$ & \\
2 & the same with $F \cap \Delta \ne \emptyset$ & makes the request on the view, with answer $v'$ and \\
  & & footprint $F'$, and writes the entry with $v'$ and $F'$; \\
  & & if $v' \ne v$, the replay \emph{resumes} \\
3 & $\mdl{n} = r$ & keeps the entry; the answer is $r$ \\
4 & $\wrt{k}{v}$ & keeps the entry and sets $k$ in the overlay \\
\bottomrule
\end{tabular}
\end{table}

Let $J$ be the journal of a run of $e$ on a state $\sigma_0$, and let $\sigma$ be another
state. A \emph{replay} of $J$ on $\sigma$ evaluates $e$ again, on the base $\sigma$, with the
journal at hand. It is given a set $\Delta$ of keys that contains every key on which
$\sigma$ and $\sigma_0$ differ. Table~\ref{tab:replay} says what each request does. The
request finds its own entry in the journal, at the position that the replay has reached. An
answer of the repository is kept if its footprint misses $\Delta$ (row 1). Otherwise the
request is made again, on the view of the replay, and the new answer and footprint replace
the old (row 2). A question to the model is not asked; the recorded answer is used (row 3).
If an answer in row 2 differs from the recorded one, the replay \emph{resumes}: the rest of
the journal is dropped, and from the next request on the evaluation is a run, with new
questions to the model. The replay returns a value, an overlay, and a new journal.

Replay costs the requests of row 2 and, after a resumption, the rest of the task. It never
costs a model call before the resumption.

\subsection{The Queue}

Let tasks $e_1, \ldots, e_m$ have runs on the same base $\sigma_0$, with overlays $\omega_i$
and journals $J_i$. A \emph{queue} takes the tasks in some order, which we write as
$1, \ldots, m$. It keeps a state, at first $s_0 = \sigma_0$. For each task in turn it replays
$J_i$ on $s_{i-1}$, with $\Delta$ the set of keys on which $s_{i-1}$ differs from $\sigma_0$,
and obtains a value, an overlay $\omega_i'$, and a journal $J_i'$. Then it sets
$s_i = s_{i-1} \upd \omega_i'$. The \emph{merged state} is $s_m$.

\section{Results}
\label{sec:results}

\subsection{A Replay Is a Run}

\begin{lemma}[Footprints]
\label{lem:footprint}
Let a read, a check, or a scan have the answer $v$ and the footprint $F$ on the view
$\sigma \upd \omega$. If $\sigma'$ agrees with $\sigma$ on $F$, the request has the answer
$v$ and the footprint $F$ on $\sigma' \upd \omega$.
\end{lemma}

\begin{proof}
For a read and a scan this is the definition of $F$. For a check, by induction on the
evaluation of its expression: the evaluation consults the base only where it reads a key that
is not in $\omega$, and each of these keys is in $F$. So the evaluation on $\sigma'$ takes
the same steps and reads the same keys.
\end{proof}

\begin{theorem}[A replay is a run]
\label{thm:replay}
Let $J_0$ be the journal of a run of $e$ on $\sigma_0$ with the value $v_0$ and the overlay
$\omega_0$. Let $\sigma$ be a state, and let $\Delta$ contain every key on which $\sigma$
differs from $\sigma_0$. Let the replay of $J_0$ on $\sigma$ with $\Delta$ return $v$,
$\omega$, and $J$. Then:
\begin{enumerate}
\item $J$ is the journal of a run of $e$ on $\sigma$ with the value $v$ and the overlay
  $\omega$. In this run the model gives the answers of $J_0$ up to the point at which the
  replay resumes.
\item If the replay does not resume, then $v = v_0$ and $\omega = \omega_0$, the journal $J$
  is $J_0$ up to footprints, and the model is not asked.
\item If $\fp{J_0} \cap \Delta = \emptyset$, the replay makes no request at all, and
  $J = J_0$.
\end{enumerate}
\end{theorem}

\begin{proof}
We show that, as long as the replay has not resumed, it has the environment and the overlay
that the run on $\sigma_0$ had at the same request, and the entries that it has written are
those of a run on $\sigma$, equal to the entries of $J_0$ up to footprints. This holds at
the start. Since the environments are equal, the next request of the replay is the next
request of $J_0$, with the same key or the same question.

\emph{Row 1.} The recorded answer $v$ with the footprint $F$ was computed on
$\sigma_0 \upd \omega$. The state $\sigma$ agrees with $\sigma_0$ outside $\Delta$, hence on
$F$. By Lemma~\ref{lem:footprint}, the answer on $\sigma \upd \omega$ is $v$ with the
footprint $F$, so the entry is the one that a run on $\sigma$ writes.

\emph{Row 2.} The entry is what a run on $\sigma$ writes, since the request was made on the
view. If the answer is the recorded one, the environments stay equal. If not, the replay
resumes, and everything after this entry is a run by definition.

\emph{Row 3.} A run may get any answer from the model, in particular $r$.

\emph{Row 4.} The key and the value are computed from the environment, which is the same.

This proves (1), and (2) is the case in which the invariant holds to the end. For (3), every
entry falls under row 1, 3 or 4.
\end{proof}

Part (1) is what makes replay more than an optimisation. The run that it constructs is one
that could have happened on $\sigma$: an agent that had started on the new state and had
received the same answers from its model would have left this journal. Up to the point of
resumption the model would have been asked the same questions, so there is no reason for it
to have answered differently.

\subsection{A Merge Is a Schedule}

\begin{theorem}[A merge is a schedule]
\label{thm:schedule}
Let a queue take tasks $e_1, \ldots, e_m$ that ran on $\sigma_0$, and let it produce the
states $s_0, \ldots, s_m$ and, for each task, the value $v_i'$, the overlay $\omega_i'$ and
the journal $J_i'$. Then for every $i$, $J_i'$ is the journal of a run of $e_i$ on $s_{i-1}$
with the value $v_i'$ and the overlay $\omega_i'$, and $s_i = s_{i-1} \upd \omega_i'$. So
the merged state $s_m$ is the state after the tasks have run one after the other in the
order of the queue, each on the state that the one before left. In these runs the model
gives each task the answers that it gave on $\sigma_0$, up to the point at which the task
resumes.
\end{theorem}

\begin{proof}
By Theorem~\ref{thm:replay}(1) for each task, with $\sigma = s_{i-1}$.
\end{proof}

The theorem holds for every order. The order decides what the queue costs.

\begin{definition}[Stale reads]
\label{def:graph}
For runs on $\sigma_0$, task $j$ \emph{changes} a key $k$ if $\omega_j(k)$ is defined and
differs from $\sigma_0(k)$. The \emph{graph of stale reads} has the tasks as its vertices and
an edge from $i$ to $j$ if $i \ne j$ and $\fp{J_i}$ contains a key that $j$ changes.
\end{definition}

\begin{corollary}[Merging for free]
\label{cor:free}
If the graph of stale reads has no cycle, take the tasks in an order in which every edge goes
from an earlier task to a later one. Then the queue makes no request, every journal is kept,
and the merged state is $\sigma_0 \upd \omega_1 \upd \cdots \upd \omega_m$.
\end{corollary}

\begin{proof}
By induction along the order. Before task $i$, the journals of the earlier tasks have been
kept, so $s_{i-1}$ differs from $\sigma_0$ only on keys that earlier tasks change. No edge
goes from $i$ to an earlier task, so $\fp{J_i}$ contains none of these keys, and
Theorem~\ref{thm:replay}(3) applies.
\end{proof}

A task that read what another wrote goes first. This is the order of a serialization
graph~\citep{fekete2005}: the reader saw the old value, so in an equivalent sequence it comes
before the writer. Two tasks may write the same key and still merge for free, if neither
read it; the later one wins, as it would in a sequence.

\begin{corollary}[Who resumes]
\label{cor:feedback}
Let $S$ be a set of tasks such that the graph of stale reads without $S$ has no cycle. Take
the other tasks first, in an order along the edges, and then the tasks of $S$. Then only
tasks of $S$ make a request or resume. For two tasks, at most one resumes, and none unless
each has read a key that the other changes.
\end{corollary}

\begin{proof}
The tasks outside $S$ merge as in Corollary~\ref{cor:free}.
\end{proof}

\begin{figure}[t]
\[
\begin{array}{@{}l@{}}
\textbf{merge}(\sigma_0;\ e_1, \ldots, e_m \text{ with journals } J_i \text{ and overlays } \omega_i)\\
\quad G := \text{the graph of stale reads of } J_1, \ldots, J_m \text{ and } \omega_1, \ldots, \omega_m\\
\quad S := \text{a set of tasks such that $G$ without $S$ has no cycle}\\
\quad \textit{order} := \text{the tasks outside $S$ along the edges of $G$, then the tasks of $S$}\\
\quad s := \sigma_0\\
\quad \textbf{for } i \textbf{ in } \textit{order}\textbf{:}\\
\qquad \Delta := \text{the keys on which $s$ differs from $\sigma_0$}\\
\qquad (v_i', \omega_i', J_i') := \text{the replay of $J_i$ on $s$ with $\Delta$, for the task $e_i$}\\
\qquad s := s \upd \omega_i'\\
\quad \textbf{return } s
\end{array}
\]
\Description{Pseudocode of a merge queue that orders the tasks by the graph of stale reads
and replays their journals.}
\caption{A queue that chooses its order. By Theorem~\ref{thm:schedule} the result is that of
the tasks run in this order, whatever the set $S$; by Corollary~\ref{cor:feedback} only tasks
of $S$ cost anything.}
\label{fig:queue}
\end{figure}

Figure~\ref{fig:queue} puts the two corollaries together. A smallest set $S$ is a minimum
feedback vertex set, which is hard to find in general~\citep{karp1972} and easy for the few
tasks that share a cycle in practice. The
corollary gives an upper bound: a task of $S$ whose footprint has changed need not resume,
since its requests may still give the recorded answers.

Cycles cannot be argued away.

\begin{proposition}[Cycles are real]
\label{prop:cycle}
For every $n \ge 2$ there are $n$ tasks and a check such that the graph of stale reads is a
cycle through all of them, every order of the tasks ends in a state in which the check
holds, and the union of their writes is a state that no order produces and in which the
check fails.
\end{proposition}

\begin{proof}
Take keys $0, \ldots, n - 1$, all $1$, and let task $i$ set key $i$ to $0$ if key
$i + 1 \bmod n$ is $1$. The check holds if some key is $1$. On the base every task writes,
so the union is all $0$. Now take any order, and suppose that it ends with all keys at $0$.
Key $i$ is set only by task $i$, so every task wrote. The last task wrote only if the key of
its neighbour was still $1$; but the neighbour ran earlier and had written. So some key is
$1$ at the end.
\end{proof}

\subsection{Checks Travel with Their Footprints}

\begin{corollary}[Checks]
\label{cor:checks}
Let the journal $J_m'$ of the last task of a queue end with entries
$\chk{c_1} = v_1, \ldots, \chk{c_n} = v_n$. Then $v_1, \ldots, v_n$ are the values of these
checks in the merged state. If the task did not resume, the queue ran, of these checks, only
those whose recorded footprints contain a key that an earlier task changed.
\end{corollary}

\begin{proof}
By Theorem~\ref{thm:schedule}, $J_m'$ is the journal of a run on $s_{m-1}$. No write follows
the checks, so they were evaluated on the final view of the task, which is the merged state.
The second claim is row 1 of Table~\ref{tab:replay}.
\end{proof}

If every task ends with the test suite of the repository, the suite holds after every merge,
and what a merge costs in tests is what the footprints say. A verdict is good for every
state that agrees with the state on which it was obtained on the keys that the test read.
Build systems call such a record a verifying trace~\citep{mokhov2018}. The corollary says
that a verifying trace is as good across a merge as it is across a rebuild.

\section{Fleets}
\label{sec:fleet}

How many agents can work on one repository at once? The theorems turn the question into one
about a random graph.

\subsection{A Model}

Let the repository have $n$ keys. Let each of $k$ tasks read $r$ keys, chosen uniformly at
random, and change $w$ of the keys that it read. The tasks run at the same time on the same
base and are then merged. The probability that a given task read a key that another given
task changed is
\[
p = 1 - \binom{n - r}{w} \Big/ \binom{n}{w} \approx \frac{r w}{n}.
\]
A task is \emph{affected} in a queue if an earlier task changed a key that it read. We count
affected tasks as tasks that resume. This is the worst case; an affected task whose
answers stand costs less.

\paragraph{The order of arrival.} If the queue takes the tasks as they come, the task in
position $i$ is affected with probability $1 - (1 - p)^{i - 1}$, since the writes of
different tasks are independent. The expected number of affected tasks is
\[
A(k) = k - \frac{1 - (1 - p)^k}{p} \approx \frac{p\,k^2}{2} \quad \text{while $p k$ is small.}
\]

\paragraph{An order along the edges.} By Corollary~\ref{cor:feedback} a queue can confine
the affected tasks to a feedback vertex set of the graph of stale reads. While $p k$ is
below one, the cycles of the graph are almost all of length two: pairs of tasks that read
each other's writes. Such a pair has probability $p_2 \approx w^2 / n + p^2$. The first term
is the chance that the two changed a common key, which each of them also read. So the
expected number is about $p_2\,k^2 / 2$. When $p k$ exceeds one, a random directed graph
with this density has a strongly connected component that contains a constant fraction of
its vertices~\citep{karp1990}, and no order helps much.

\subsection{Numbers}

\begin{table}[t]
\caption{A fleet of $k$ tasks on $10{,}000$ keys; each task reads $50$ keys and changes $5$
of them, so that $p$ is \exFleetP{} percent. The expected number of affected tasks in the
order of arrival is the closed form $A(k)$. For an order along the edges it is the mean of
$300$ samples, with the feedback set found greedily, so it is an upper bound. The last
column is the share of samples whose graph has a cycle.}
\label{tab:fleet}
\small
\begin{tabular}{@{}rrrrr@{}}
\toprule
$k$ & $p\,k$ & order of arrival & order along the edges & with a cycle \\
\midrule
10 & \exFleetKpa & \exFleetArrivea & \exFleetBesta & \exFleetCyclica\% \\
20 & \exFleetKpb & \exFleetArriveb & \exFleetBestb & \exFleetCyclicb\% \\
40 & \exFleetKpc & \exFleetArrivec & \exFleetBestc & \exFleetCyclicc\% \\
80 & \exFleetKpd & \exFleetArrived & \exFleetBestd & \exFleetCyclicd\% \\
160 & \exFleetKpe & \exFleetArrivee & \exFleetBeste & \exFleetCyclice\% \\
\bottomrule
\end{tabular}
\end{table}

Table~\ref{tab:fleet} shows a repository of $10{,}000$ keys in which a task reads $50$ and
changes $5$. With $40$ tasks at once, $p k$ is about one. In the order of arrival
\exFleetArrivec{} of the $40$ tasks are affected. In an order along the edges it is
\exFleetBestc{}. Here $p_2$ is \exFleetMutual{} percent in a sample of $200{,}000$ pairs, so $p$ is \exFleetRatio{}
times as large, and most of $p_2$ is the term
$w^2 / n$: the pairs that cannot be ordered are mostly pairs that edited the same key. With
$160$ tasks the graph is one large tangle, and a third of the fleet is affected in the best
order that we find.

\paragraph{The size of a fleet.} Suppose an affected task costs a fraction $\rho$ of a
task, and the affected tasks are done again one after the other. A batch of $k$ tasks then
takes time $1 + \rho\,q\,k^2 / 2$ instead of $1$, where $q$ is $p$ or $p_2$, and the number
of tasks finished per unit of time is largest at
\[
k^\ast = \sqrt{\frac{2}{\rho\,q}}.
\]
At this size the fleet runs at half its nominal speed. For $\rho = 1 / 2$ and the numbers
above, $k^\ast$ is \exFleetSizeArrive{} in the order of arrival and \exFleetSizeEdges{} along
the edges. The square root
says that a fleet does not scale by adding agents: four times the agents need a repository
with a sixteenth of the overlap. Replay lowers $\rho$, and the order lowers $q$.

\subsection{A Registry That Everyone Reads}

The model is crude. Real tasks do not read uniformly: some keys are read by everyone. A
lock file, a table of routes, a list of exports, a changelog are of this kind, and so is any
file that a tool reads on every run. Add one such key to the model. Every task reads it, and
a task changes it with probability $h$.

Two effects follow. The tasks that change the registry have all read it, so each pair of
them forms a cycle, and of $m$ such tasks $m - 1$ are affected in every order. And a task
that does not change the registry has still read it, so it must go before every task that
does. If one of those has read any of its writes, the two form a cycle that would have been
an edge without the registry. With $h = 0.1$ and the numbers of Table~\ref{tab:fleet}, the
affected tasks in an order along the edges rise from \exFleetBesta{} to \exFleetHota{} for
ten tasks, from \exFleetBestb{} to \exFleetHotb{} for twenty, and from \exFleetBestc{} to
\exFleetHotc{} for forty. One key that everyone reads and a few write costs more than all
the other overlap together.

The remedy is in the choice of keys. If each entry of the registry is a key of its own, a
task that adds an entry changes a key that nobody has read, and a task that looks up an
entry reads one key. What remains expensive is a request that reads the registry as a
whole, which is $\scan$. The same holds for a directory: creating a file is cheap, and
listing the directory reads everything in it.

The next section looks at footprints in a real repository.

\section{A Library of Proofs}
\label{sec:library}

A library for a proof assistant is a repository whose checks are exact: a module either
elaborates and passes the kernel or it does not. The key question for a merge is what the
check of a module reads. We measured this on Lean's own sources, version 4.36: the
\exModules{} modules under \texttt{Init}, \texttt{Std} and \texttt{Lean}.

\subsection{Three Footprints}

A Lean module has a \emph{public scope}, which its importers see, and a \emph{private scope}:
the bodies of proofs, of definitions that are not exposed, and whatever is declared
private~\citep{leanmodules}. An import brings in the public scope of the imported module. A
\texttt{public import} passes it on to further importers, and an \texttt{import all} brings
in the private scope as well. The library has \exImportsPublic{} public imports,
\exImportsPlain{} plain ones, and \exImportsAll{} imports of a private scope. We compare
three accounts of what the check of a module $M$ reads.
\begin{itemize}
\item \emph{By files.} Everything that $M$ imports, directly or through other modules. This
  is the footprint without a module system, where any change of a file makes all its
  importers be built again.
\item \emph{Public scopes.} The public scopes that $M$ sees: those of its imports, of what
  they import publicly, and so on, and the same for the imports of modules whose private
  scope $M$ sees.
\item \emph{Private scopes.} The private scopes that $M$ sees: those of the modules that it
  reaches through \texttt{import all}.
\end{itemize}
For a change of a module $B$, the checks that must be made again are those of the modules
whose footprint contains $B$: the footprint by files for any change of the file, the public
scopes for a change of the public scope of $B$, and the private scopes for a change that
stays in its private scope, such as a new proof of an old statement. The count for a public change is a lower bound. A module
that is checked again may come out with a different public scope, and then its importers
follow; a build system finds that out as it goes.

\subsection{Numbers}

\begin{table}[t]
\caption{The number of other modules whose check must be made again after a change of one
module of Lean's library, over all \exModules{} modules, and the share of ordered pairs of
modules in which the first must be checked again after a change of the second.}
\label{tab:cones}
\small
\begin{tabular}{@{}lrrrr@{}}
\toprule
What changed & mean & median & largest & share of pairs \\
\midrule
the file & \exConeFileMean & \exConeFileMedian & \exConeFileMax & \exConeFileP\% \\
the public scope (at least) & \exConePublicMean & \exConePublicMedian & \exConePublicMax & \exConePublicP\% \\
the private scope only & \exConePrivateMean & \exConePrivateMedian & \exConePrivateMax & \exConePrivateP\% \\
\bottomrule
\end{tabular}
\end{table}

Table~\ref{tab:cones} shows the result. By files, a change of a module makes
\exConeFileMean{} modules be checked again on average, a fifth of the library, and for the
modules at the bottom nearly all of it: \exConeFileHalf{} modules are imported, directly
or not, by at least half of the library. For a change of a public scope the mean is
\exConePublicMean{}, and \exConePublicHalf{} modules are seen by half of the library. For a
change that stays in a private scope the mean is \exConePrivateMean{} and the largest count
is \exConePrivateMax{}; for \exConePrivateNone{} percent of the modules no other module
sees the private scope at all.

So two agents that both change only proofs in this library almost never invalidate each
other's checks, and by Corollary~\ref{cor:checks} their work merges with nothing checked
again, apart from the few modules that reach into private scopes. Two agents of whom one
changes a statement are in the other regime: hundreds of checks are made again. Whether any
agent has to resume is then decided by what those checks say. An earlier paper in this
series measured the same boundary from the side of a human reader, who must read statements
and may skip proofs~\citep{tpp1}. The reader and the merge queue profit from the same fact:
a proof is read by nobody.

\subsection{What Does Not Stay Local}

The measurement is about modules. Inside the elaboration of a module there are reads that no
scope bounds, and they behave like $\scan$. Resolving a type class instance consults every
instance in scope. A simplifier call consults every lemma that is marked for it. Notations
and coercions are looked up in tables to which any imported module may have added. A new
instance in an imported module can change what a term elaborates to, even if no declaration
that the term mentions has changed. Our counts treat the public scope of a module as one
key, which is the cautious choice: any change of it counts. A finer account, with
declarations as keys, would have to treat these tables as keys that everybody reads.

The kernel is better placed than the elaborator. A proof term that has been elaborated
mentions the constants that it uses, and its check reads their statements, and their bodies
where it unfolds them. In a system with global universe constraints the kernel also reads a
table that every declaration may extend; Lean's universe polymorphism has no such table. This
is the pattern of the whole section: what is looked up by name is local, and what is looked
up by search is not.

\section{Testing the Claims}
\label{sec:tests}

The artifact is one Python script without dependencies, and a second one that extracts the
import graph and needs a Lean toolchain; the graph is committed, so
\texttt{make -C artifact check} needs only Python. All numbers are those of one fixed seed.

\paragraph{Tasks.} A repository has six keys with values below four, and three checks. Tasks
and checks are random expressions of the calculus, of nesting depth six and three. A model
is a random stream of numbers.

\paragraph{Theorem~\ref{thm:replay}.} For each of \exRebases{} tasks the script runs the
task, gives each key of the base a new random value with probability $0.3$, and replays the
journal on the new state. It then runs the task live on the new state, with a model that gives the answers
of the new journal, and checks that the live run has the same journal, with the same
footprints, the same overlay, and the same value. It also checks parts (2) and (3). Of the
replays, \exRebaseResumed{} resume. In those, \exResumedKeptPercent{} percent of the
model's answers lie before the point of resumption and are kept.

\paragraph{Theorem~\ref{thm:schedule}.} For each of \exQueues{} sets of two to four tasks
the script runs the tasks on one base, merges them in a random order, and compares the
queue with the tasks run one after the other, live, with the model answers of the new
journals: every journal, every overlay, and the final state. Of the \exQueueTasks{}
replays, \exUntouched{} make no request, \exCompared{} make requests again and find the old
answers, and \exResumed{} resume. The queues keep \exQueueSaved{} answers of the model and
ask for \exQueueCalls{}, so \exSavedPercent{} percent of the calls are saved against running
every task again. Of \exEntries{} recorded reads, checks and scans, \exRepeated{} are made
again, which is \exRepeatedPercent{} percent.

\paragraph{The corollaries.} In \exFree{} of the sets the graph of stale reads has no cycle.
The script merges them in an order along the edges and checks that no request is made and
that the merged state is the union of the writes (Corollary~\ref{cor:free}). In the other
\exCyclic{} it computes a smallest feedback vertex set by search and checks that only its
tasks make requests (Corollary~\ref{cor:feedback}). For Corollary~\ref{cor:checks} it runs
$1{,}000$ queues of tasks that end with the three checks and compares the verdicts in the
last journal with the checks evaluated on the merged state. For Proposition~\ref{prop:cycle}
it builds the rings of two, three and four tasks and tries every order.

\begin{table}[t]
\caption{Five queues with one weakened step. Each ends, for some set of tasks, in a state
that is not the state of the tasks run one after the other. The last column is the number of
random sets that were tried until the first such set.}
\label{tab:broken}
\small
\begin{tabular}{@{}llr@{}}
\toprule
Step & What the weakened queue does & Tried \\
\midrule
row 2, reads & keeps the answer of a read whose key has changed & \exStepSkipRead \\
row 2, checks & keeps the verdict of a check whose footprint has changed & \exStepSkipCheck \\
footprints & gives $\scan$ an empty footprint & \exStepScanLocal \\
$\Delta$ & takes the writes that the earlier tasks made on the base & \exStepStaleDelta \\
$\Delta$ & takes the writes of the task just before & \exStepLastDelta \\
\bottomrule
\end{tabular}
\end{table}

\paragraph{Every step is needed.} Table~\ref{tab:broken} weakens the replay in five ways.
For each the script finds a set of tasks for which the weakened queue ends in a state that
differs from the tasks run one after the other. The fourth row is the least obvious. A queue
that computes $\Delta$ from the writes that the earlier tasks made on the base forgets that
a task which resumed may have written other keys.

\paragraph{Fleets and the library.} The script checks the closed form $A(k)$ against the
samples of Table~\ref{tab:fleet} and computes Table~\ref{tab:cones} from the committed
graph.

\paragraph{What the tests do not show.} The tasks are small and the repository has six keys.
The fleet is a model with uniform footprints. The numbers of Section~\ref{sec:library} rest
on our reading of what a check reads, by modules, and on the import lines of the sources; we
did not build the library and watch which modules were built again.

\section{Limits}
\label{sec:limits}

\paragraph{The journal must be complete.} The theorems assume that everything that a task
learns about the repository goes through requests that are journaled. An agent that reads a
file through a side door has a footprint that the journal does not know. The companion
paper on journals gives the discipline that excludes this~\citep{tpp8}.

\paragraph{Keys.} A key is the unit of change and of footprint. With files as keys, two
edits of different functions in one file are a conflict. Smaller keys give smaller
footprints and need a way to tell which declaration a tool read.

\paragraph{The model's answers.} Replay reuses the answers of the model before the point of
resumption. This is justified in the sense of Theorem~\ref{thm:replay}: the questions were
the same. It is not justified for a model that looks at the repository on its own, outside
the requests of the task.

\paragraph{Intent.} A merged state that is the outcome of a schedule is correct in the
sense that each task did, on the state it found, what it would have done there. Whether two
tasks make sense together is another question. Two agents can each add a feature that works
and that the user would not want side by side. Only checks express such requirements, and
only if someone wrote them.

\paragraph{Creation and deletion.} The calculus has a fixed set of keys. A task that lists a
directory reads the set of names, and a task that creates a file changes it. This fits the
model if the set of names is treated as a key that every listing reads, which makes listings
expensive, as $\scan$ is.

\paragraph{Proofs.} The proofs are on paper, and the artifact tests them and does not prove
them.

\section{Related Work}
\label{sec:related}

\paragraph{Isolation and serializability.} That a history of transactions is correct if it
is equivalent to a serial one is due to \citet{papadimitriou1979}. Snapshot isolation and
its anomaly, write skew, were described by \citet{berenson1995}; \citet{fekete2005}
characterised when snapshot isolation is serializable by cycles in a graph of dependencies,
and \citet{cahill2008} turned the characterisation into a run-time mechanism. Our graph of
stale reads is this graph for tasks that commit together, and Corollary~\ref{cor:free} is
the serializable case. Optimistic concurrency control~\citep{kung1981} lets transactions run
without locks and validates their reads before they commit; a transaction that fails
validation starts again. Software transactional memory does the same with a log of
reads~\citep{harris2005}. What we add is the journal: a task that fails validation does not
start again. It resumes at the first answer that differs, and keeps the expensive answers
before it.

\paragraph{Invariants under merging.} \citet{bailis2014} ask when concurrent operations can
be merged without coordination such that the invariants of an application are preserved, and
answer with a confluence condition on invariants. Corollary~\ref{cor:checks} is a dynamic
form of such a condition: an invariant that was checked survives a merge if the merge stays
outside the footprint of its check. Replicated data types merge without coordination by
making operations commute~\citep{shapiro2011}; the edits of agents do not commute in
general, and the footprints say for which pairs it does not matter.

\paragraph{Build systems and incremental computation.} A build system skips a task whose
inputs have not changed, and the systems with dynamic dependencies record what a task read
while it ran~\citep{mokhov2018}. Self-adjusting computation records the reads of a whole
program and, after a change of its input, runs again the parts that read what
changed~\citep{acar2002}. Replay is change propagation for a program with an oracle in it:
the parts before the first affected read are kept with their answers of the oracle, and the
rest is run again because the oracle may now be asked something else.

\paragraph{Version control.} Theories of patches explain when two edits commute and what a
merge is as a universal construction~\citep{mimram2013,angiuli2014}. They work on the text of
the edits. That a merge can be clean as text and wrong as a program is documented for real
projects~\citep{brun2011}. Our contribution is on the side of the reads, which version
control does not see.

\paragraph{Agents and libraries.} The journal and its replay come from durable execution,
where they serve recovery from crashes~\citep{tpp8}; here the same journal serves merging.
What a call of the model costs, and so what a resumption wastes, is the subject
of~\citet{tpp6}. That statements are what others read and proofs are not is the measurement
of~\citet{tpp1}, taken there for human readers and here for checks.

\section{Conclusion}

A fleet of agents on one repository is a set of transactions that commit together, and the
journals of the agents say exactly what each transaction read. With them, merging stops
being an act of faith followed by a full run of the tests. A queue can tell which tasks
merge for nothing, which need a few tests run again, and which must go back to work, and for
the last kind it can put them back at the request where their world changed. How large a
fleet can be is then a question about how much the tasks of the fleet read of what the
others write, and for a library of proofs with scopes the answer is that those who write
proofs hardly read each other at all.

\bibliographystyle{ACM-Reference-Format}
\bibliography{refs}

\end{document}
